\chapter{High-Level and Low-Level Design}
\label{ch:design}

This chapter maps the requirements and methodology to a service-oriented prototype design that can be implemented incrementally on the team's available machines. The decomposition protects experimental determinism by assigning the simulation clock and all signal actuation to one adapter. The API, dashboard, perception worker, controllers, routing logic, and experiment store exchange typed messages without competing for that authority. The design also separates a requested decision from the action applied after safety checks, preserves the source and effective time of every traffic snapshot, and binds corridor requests to an immutable route version. These choices support replay and fault diagnosis without requiring a prematurely distributed deployment. Interfaces may initially be implemented as in-process modules, provided that they retain the same ownership boundaries and contracts. The chapter moves from architecture and domain objects to interaction sequences, state transitions, data contracts, synchronization, failure handling, module policies, and testing so each important requirement has a visible design consequence.

\section{System Overview}

The system must support repeatable research runs, an interactive demonstration, and post-run analysis without allowing presentation concerns to alter experimental state. Its main design goals are deterministic ownership of state changes, replaceable controller and perception implementations, explicit time and units, safe corridor transitions, immutable run configuration, and visible failure. Performance is optimized only after these properties are preserved because a faster result is not useful when its source, phase legality, or causal order is uncertain. The design assumes one machine or a small development deployment for Deliverable 2 and FYP-I. Logical separation therefore does not imply a network service for every box in the diagrams. The orchestrator, controllers, routing logic, and repositories may begin as in-process modules behind typed interfaces, while video inference can move to a bounded worker when profiling shows it would block the control loop. This approach keeps deployment manageable but preserves a path to isolate resource-heavy components without changing experimental semantics.

\subsection{Operator Interface and Commands}

The operator uses a React dashboard for configuration, monitoring, emergency, and analysis tasks. Commands pass through FastAPI validation before reaching the experiment service or orchestrator. The interface cannot step SUMO, apply a phase, or alter an active run's immutable configuration. Responses include source, simulation step, freshness, and status so the dashboard can distinguish current evidence from a stale or unavailable view. Configuration commands resolve network, controller, seed, duration, incident, mission, and metric identifiers before a run begins. Monitoring queries expose requested and applied actions separately, while emergency commands return the route version and corridor state affected by the request. The evaluator can inspect completed and failed runs without receiving mutation privileges. A disconnected browser may reconnect from the current snapshot and event cursor, but it cannot replay an old command automatically. These rules make the graphical interface useful for demonstration and review without creating a second, poorly controlled path into the traffic simulation.

\subsection{Traffic Observation and Simulation}

SUMO provides authoritative traffic state for controlled experiments, while a bounded perception worker supports recorded-video observation. The traffic-state service normalizes either source without presenting video proxies as simulator measurements. Only the SUMO/TraCI adapter owns simulation stepping, subscriptions, and applied signal changes. Perception reports model version, input provenance, confidence, and freshness, preserving a clear record of what information was available to each decision. A simulator snapshot can include exact vehicle, lane, signal, and route values available through TraCI. A video snapshot contains only quantities justified by the selected detector, tracker, road regions, and calibration; missing physical values remain unavailable instead of being replaced by zero. Both snapshots identify the time at which values became effective and the period for which they remain valid. The control loop uses one declared source mode per experiment, while offline perception jobs can be assessed independently against labelled footage. This boundary prevents the report from treating visual detection output as equivalent to controlled simulator ground truth.

\subsection{Adaptive Control and Emergency Routing}

Normal traffic control and emergency priority share the safety filter but remain separate decision concerns. Fixed-time, pressure-inspired, and learned controllers return requested actions through a common interface. Routing creates immutable, time-costed ambulance paths, and the corridor coordinator translates the current route into bounded priority requests. Neither module actuates a signal directly; every request must satisfy legal phase, clearance, ownership, expiry, and recovery rules. The common controller interface supplies the same snapshot, legal phases, decision interval, and explanation fields to each normal strategy so comparison is not distorted by different privileges. Emergency operation is supervisory: it temporarily owns selected intersections through an explicit state machine and returns them to a valid normal boundary after passage, cancellation, or timeout. Rerouting creates a new route version and cancels obsolete future requests before activating the revised look-ahead set. Separating route choice, corridor priority, normal control, and safety enforcement permits useful ablation experiments and makes a rejected decision traceable to its responsible component.

\subsection{Experiment and Evaluator Interaction}
\label{sec:architecture}

\begin{figure}[htbp]
\centering
\begin{tikzpicture}[font=\footnotesize,node distance=7mm and 11mm,
comp/.style={draw,rounded corners,fill=manzilgray,minimum width=3.35cm,minimum height=1.0cm,align=center},
ext/.style={draw,rounded corners,fill=manzilblue!12,minimum width=3.35cm,minimum height=1cm,align=center},
arr/.style={-{Latex},thick}]
\node[ext] (ui) {React dashboard\\operator and evaluator};
\node[comp,below=of ui] (api) {FastAPI boundary\\validation and queries};
\node[comp,below=of api] (orch) {Control orchestrator\\ordering and ownership};
\node[comp,left=of orch] (exp) {Experiment service\\config, lifecycle, metrics};
\node[comp,right=of orch] (state) {Traffic-state service\\typed snapshots};
\node[comp,below=of exp] (store) {Experiment store\\repository interface};
\node[comp,below left=of orch] (controller) {Controller registry\\fixed / rule / RL};
\node[comp,below=of orch,fill=manzilgreen!12] (safety) {Safety filter\\legal phase transitions};
\node[comp,below right=of orch] (route) {Routing and corridor\\mission lifecycle};
\node[comp,below=of safety,fill=manzilgreen!12] (sumoad) {SUMO/TraCI adapter\\sole step and actuation owner};
\node[ext,below=of sumoad] (sumo) {SUMO process};
\node[comp,below=of state] (per) {Perception worker\\detector / tracker};
\node[ext,below=of per] (video) {Recorded video};
\draw[arr,<->] (ui)--node[right]{REST / stream}(api);
\draw[arr] (api)--(exp);\draw[arr] (api)--(orch);\draw[arr] (state)--(api);
\draw[arr,<->] (orch)--(controller);\draw[arr,<->] (orch)--(route);\draw[arr] (orch)--(safety);\draw[arr] (safety)--(sumoad);
\draw[arr,<->] (sumoad)--(sumo);\draw[arr] (sumoad)--(state);
\draw[arr] (video)--(per);\draw[arr] (per)--(state);
\draw[arr,<->] (exp)--(store);\draw[arr] (orch)--(exp);\draw[arr] (state)--(exp);
\end{tikzpicture}
\caption{Proposed high-level architecture}
\label{fig:architecture}
\end{figure}

Figure \ref{fig:architecture} shows the command and evidence paths around the single-owner simulation core. The React dashboard sends configuration or mission commands through FastAPI and receives non-authoritative views; it never calls TraCI. The orchestrator follows Algorithm \ref{alg:loop}, requests decisions from the selected normal controller or corridor coordinator, and sends validated actions to the adapter. Only that adapter advances SUMO and applies signal changes. The traffic-state service constructs source-specific snapshots from simulator observations or perception messages while preserving effective time, units, confidence, and validity. Routing maintains mission and route versions, and the corridor coordinator owns intersection-level priority lifecycles. The experiment service validates run configuration, records events and metrics, and persists them through a repository interface whose proposed implementation is PostgreSQL. Directional arrows in the figure distinguish state flow from command ownership, making clear that perception can inform a snapshot without acquiring authority over simulation time. This separation also keeps dashboard latency outside the control path.


Table \ref{tab:components} assigns each component the data it consumes and the state or evidence it is permitted to produce. This allocation prevents shared mutable state from becoming an undocumented integration mechanism. The dashboard owns only presentation state, controllers own only proposed actions and explanations, and the SUMO adapter owns the external process, subscriptions, clock, and applied actuation. The orchestrator orders these modules but does not replace their domain logic. Routing and corridor responsibilities are listed separately because a shortest path does not imply authority to change a signal, and a corridor request must retain the route version that justified it. The traffic-state service can normalize observations from two source modes but may not silently combine missing video measures with exact simulator values. Finally, the experiment store receives validated records and artifacts rather than arbitrary objects from every module. These boundaries give each requirement a testable owner and permit an implementation to remain modular even when several components share one process.

\begin{longtable}{p{0.23\textwidth}p{0.30\textwidth}p{0.36\textwidth}}
\caption{Component responsibility allocation}\label{tab:components}\\
\toprule
\textbf{Component} & \textbf{Consumes} & \textbf{Produces / owns}\\
\midrule\endfirsthead
\toprule
\textbf{Component} & \textbf{Consumes} & \textbf{Produces / owns}\\
\midrule\endhead
React dashboard & Query responses, event snapshots, user input & Presentation state; no control authority\\
FastAPI boundary & HTTP/stream commands, identity-free local operator session & Validation, command/query boundary, error mapping\\
Experiment service & Valid configuration, metric/event batches & Run lifecycle, immutable manifest, comparison eligibility\\
Control orchestrator & Run config and ordered state snapshot & Decision schedule, module ordering, requested actions\\
SUMO/TraCI adapter & Validated step and action commands & Simulation process, clock, subscriptions, applied actions\\
Traffic-state service & SUMO observations or perception messages & Versioned source-specific snapshots and freshness\\
Perception worker & Video and model configuration & Detections, tracks, aggregates, job health\\
Controller registry & State, legal phases, timing constraints & Requested normal-control action and explanation\\
Routing service & Graph, mission, cost snapshot & Route versions, no-route reason, reroute decision\\
Corridor coordinator & Route version, ambulance state, signal state & Priority request/state, release and recovery requests\\
Experiment store & Validated records and artifacts & Durable configuration, events, metrics, checksums\\
\bottomrule
\end{longtable}

\section{Design Considerations}

\subsection{Assumptions and Dependencies}

The design assumes a validated SUMO network, legal traffic-light programs, reproducible demand files, and approved recorded video when perception experiments are performed. It depends on TraCI for simulator control and on versioned model artifacts and repositories for repeatable evidence. A single development machine is sufficient for initial integration, although detector training may require separate GPU capacity. Missing maps, checkpoints, or video sources therefore produce explicit fallbacks or blocked experiments rather than hidden substitutions. Network preparation must resolve disconnected edges, prohibited movements, signal programs, and ambulance-compatible routes before a scenario is accepted. Demand generation exposes its seed and distribution parameters so repeated runs can be paired. Video processing assumes that clip provenance, model checksum, road regions, and any calibration are available. The repository must retain configuration and causal events even when large artifacts remain external files. These assumptions are verified at configuration time where possible; unmet conditions become validation errors, declared limitations, or a documented synthetic-network fallback rather than undocumented changes to the experiment.

\subsection{General Constraints}

The prototype cannot assume city-wide live sensors, traffic-authority deployment, or emergency-service integration. Evaluation is simulation based, with recorded video used only for perception validation. The core scope permits one active emergency mission, every signal transition must follow the validated program, and a run's configuration becomes immutable at start. Failed runs remain visible, while no dashboard action can bypass backend validation. These constraints prevent a polished demonstration from being confused with operational deployment. The selected Lahore network must remain small enough for repeated seeded experiments and readable inspection of imported connections. The system does not claim certified preemption, field latency, legal authority, or direct control of real signals. Learned controllers operate within the same legal action space and decision schedule as baselines, and training-only tuning cannot use held-out evaluation scenarios. Resource limits may require offline detector inference, compact networks, or a documented synthetic grid. Such reductions are acceptable only when the report discloses them and conclusions remain bounded to the tested network, demand, footage, models, and computing environment.

\subsection{Goals and Guidelines}

The design prioritizes traceability, deterministic ordering, safe actuation, replaceable research components, and honest reporting of unavailable evidence. Each mutable resource has one owner; messages carry version and effective time; requested and applied actions remain distinct; and fallback behavior is bounded and labelled. Interfaces are small enough for contract testing, and diagrams describe logical responsibilities rather than claiming every component is already deployed as a separate service. Naming, units, source modes, and terminal states remain consistent across the API, database, logs, interface, and report. A component rejects an unknown schema, expired snapshot, illegal phase, disconnected route, or incompatible comparison rather than trying to infer intent. Experiments retain failures, exclusions, and fallback intervals so favorable summaries cannot hide them. Performance optimization follows profiling and must not bypass validation or remove evidence. A change is accepted only when its owner, inputs, outputs, time semantics, failure behavior, and effect on reproducibility can be stated and tested.

\subsection{Development Methods}

Table \ref{tab:decisions} records the main alternatives considered and the consequence of each selected approach. The choices favor traceability and feasible integration over maximum distribution or algorithmic novelty. SUMO remains the simulator because it aligns with the approved proposal and provides detailed route and signal inspection, although CityFlow could offer greater training throughput for a different scope. A common controller interface and safety filter make fixed, rule-based, and learned policies comparable, while supervisory emergency control permits a separate corridor ablation. Dijkstra remains the routing reference and A* is conditional on an admissible time lower bound. Structured persistence is proposed because flat files make cross-run relationships difficult to enforce, but the repository boundary keeps the engine replaceable. The decision record is not evidence that every selected technology has been completed; it explains why the design adopts these constraints and what cost each creates, including additional message schemas, job states, and performance profiling.

\begin{longtable}{P{0.20\textwidth}P{0.34\textwidth}P{0.35\textwidth}}
\caption{Architectural decision record}\label{tab:decisions}\\
\toprule
\textbf{Decision} & \textbf{Selected approach and consequence} & \textbf{Alternative considered}\\
\midrule\endfirsthead
\toprule
\textbf{Decision} & \textbf{Selected approach and consequence} & \textbf{Alternative considered}\\
\midrule\endhead
Simulation ownership & One orchestrator and adapter control stepping and signals, giving a deterministic replay and concurrency model. & Multiple services calling TraCI risk non-deterministic order.\\
Simulator & SUMO aligns with the proposal and supplies detailed road and signal tooling; performance is profiled later. & CityFlow offers training throughput but changes the approved environment.\\
API and interface & FastAPI and React provide a typed boundary and independent interface iteration. & A server-rendered monolith is simpler to deploy but weakens service contracts.\\
Controller integration & A common request interface and safety filter make controllers comparable and preserve applied-action evidence. & Direct TraCI access couples research policies to actuation.\\
Emergency control & A supervisory corridor state machine supports a clear ablation and one safety owner. & Folding emergency reward into normal RL confounds attribution.\\
Routing & Dijkstra is the reference; A* is permitted only with an admissible time lower bound. & Distance-only routing does not satisfy the congestion-aware evaluation objective.\\
Persistence & A structured experiment schema remains portable behind a repository interface. & Flat files are quick but weak for cross-run integrity.\\
Perception execution & A bounded background worker prevents inference from blocking control while exposing job state. & In-request inference couples API latency to model runtime.\\
\bottomrule
\end{longtable}

\section{System Architecture}

The Manzil architecture follows the component boundaries introduced in Section \ref{sec:architecture} and refines them into objects that can be implemented and tested independently. The dashboard and API form the presentation boundary, while the orchestrator, traffic-state service, controller registry, routing service, corridor coordinator, safety filter, simulation adapter, and experiment store form the application core. This partition keeps visualization outside the authoritative control path and assigns every mutable resource to one owner. Communication may initially use in-process calls, but each interaction retains a typed request or event contract so that later separation does not change experimental meaning. The architecture also distinguishes a proposed action from an applied action, a traffic observation from its source, and a route from the version that authorized corridor activity. These distinctions make replay and fault diagnosis possible because each state change can be connected to its input snapshot, policy, safety decision, and effective simulation step. The following domain diagram records these relationships without committing the implementation to one framework or database library.

\subsection{High-Level Partitioning}

The presentation, application, research, infrastructure, and evidence responsibilities form the principal partitions. React and FastAPI form the presentation boundary; the orchestrator and traffic-state service coordinate the application flow; and controllers, perception, routing, and corridor logic contain replaceable research behavior. The SUMO adapter and repositories isolate external infrastructure, while the experiment service connects every partition to durable evidence. Figure \ref{fig:architecture} shows their communication and Table \ref{tab:components} states the data each partition may own. Partition boundaries are logical rather than mandatory network processes. During initial integration, several application modules may run inside one Python process, provided that they exchange the same validated objects and respect ownership. The perception worker may be isolated when inference affects timing, and the dashboard remains a separate client because its refresh rate must not control simulation progress. This partitioning provides a manageable deployment for the team while retaining clear seams for profiling, replacement, testing, and later distribution.

\subsection{Component Collaboration}

Collaboration follows typed requests and events instead of shared mutable objects. At each step the adapter supplies state, the traffic-state service creates a snapshot, the selected controller proposes an action, the safety filter validates it, and the adapter applies the result. Routing publishes immutable route versions to the corridor coordinator. The experiment service records configuration, actions, mission changes, and metrics. Section \ref{sec:sequences} makes this ordering explicit for normal control, dispatch, and rerouting. Every request carries identifiers that link it to a run, simulation step, controller or mission, and schema version. The receiving component either produces a validated response or a structured rejection; it cannot silently repair an incompatible message. Events are delivered to persistence in causal order, while dashboard streams may omit intermediate views without changing the authoritative record. These collaboration rules make integration failures observable and allow stubs to replace unfinished modules during contract testing without bypassing the same validation applied to completed components.

\subsection{Rationale for Decomposition}

The decomposition isolates research algorithms, external tools, and presentation concerns. A controller or detector can change without receiving direct TraCI authority, while a storage engine or interface library can change without altering the meaning of a snapshot or applied action. The split also permits bounded perception execution and deterministic simulation. These benefits justify explicit contracts because reproducibility and fault diagnosis are central evaluation requirements. A monolithic script could produce an early visual demonstration, but it would make phase ownership, failure containment, and comparative evaluation difficult to verify. Conversely, deploying every box as an independent service would add network failure and configuration overhead before profiling shows a need. The selected middle position preserves conceptual services behind interfaces while permitting in-process execution. It supports incremental development: fixed control and simulation can operate before perception or learning is ready, and later modules can be added without rewriting the control loop or database meaning.

\subsection{System Architecture Diagram}

Figure \ref{fig:architecture} presents the proposed architecture. Its central feature is the single actuation path from orchestrator through safety filter and SUMO adapter. State and evidence travel outward to the traffic-state and experiment services, while the dashboard remains a non-authoritative observer and command source. This direction prevents interface refreshes, detector delays, or database queries from becoming undocumented signal commands. The figure also separates routing from corridor ownership: a route is computed from a versioned cost snapshot, but only the corridor coordinator may request temporary priority along that route. Perception contributes observations through the traffic-state service and cannot call a controller or adapter directly. The experiment service receives lifecycle and evidence records but does not determine a phase. Readers should therefore interpret arrows as permitted command or evidence paths, not merely as data availability. Missing arrows are intentional constraints that preserve determinism and traceability.

\subsection{Subsystem Architecture}

Each subsystem exposes a narrow testable interface. Controllers return proposed actions, routing returns versioned paths, the corridor coordinator returns route-bound priority requests, and the traffic-state service returns source-specific snapshots. Infrastructure adapters translate these domain objects to TraCI, HTTP, streams, and persistent records. The domain model in Section \ref{sec:domain-model} describes the stable objects shared across those interfaces. The simulation subsystem owns process launch, stepping, subscriptions, and applied commands. The experiment subsystem owns configuration validation, lifecycle status, manifests, metrics, and comparison eligibility. The perception subsystem owns jobs, model provenance, detections, tracks, and derived region observations. The emergency subsystem owns missions, route versions, priority requests, passage, and recovery. These responsibilities prevent two modules from writing the same state. Each subsystem can be exercised with fixtures representing adjacent contracts, allowing tests to distinguish local logic errors from integration or external-tool failures.

\section{Architectural Strategies}

\subsection{Choice of Technologies}

Python, SUMO, TraCI, FastAPI, React, and structured persistence match the approved proposal and divide research computation from presentation. Python supports controllers, routing, perception, and analysis; SUMO supplies reproducible traffic dynamics; FastAPI defines validated contracts; and React supports independent dashboard views. Versions are written to run manifests. Final detector, tracker, reinforcement-learning library, and database-engine choices remain configuration decisions until implementation evidence confirms them. Selection also reflects team feasibility: these technologies run locally, expose documented interfaces, and support automated testing without requiring proprietary traffic infrastructure. SUMO is retained over a throughput-oriented alternative because its network, vehicle, signal, and TraCI tooling align with the approved scope. FastAPI and React add a service boundary but allow the research core and interface to evolve independently. A technology is not accepted solely because it is popular; its version, configuration, licence, input assumptions, and effect on timing or metrics must be recorded for the experiment that uses it.

\subsection{Reuse of Existing Software Components}

Manzil reuses SUMO traffic models, TraCI primitives, OpenStreetMap conversion tools, computer-vision libraries, and standard graph algorithms. Adapters bound this reuse so an external package cannot define project semantics implicitly. Imported networks are checked for connectivity and signal legality, model checkpoints retain provenance, and route costs state their units and observation time. Manzil contributes experiment, safety, corridor, audit, and comparison logic that combines these components consistently. Reuse reduces implementation risk but does not transfer responsibility for validation to a library. The network importer must still confirm turn connections and traffic-light programs; the perception pipeline must still evaluate selected checkpoints and tracking rules on the intended footage; and the routing service must still build a legal vehicle-movement graph. Third-party defaults are captured in manifests when they affect results. External components are accessed through project-owned interfaces so a library upgrade or replacement produces an explicit compatibility change instead of silently altering stored evidence or control behavior.

\subsection{Future Extension and Enhancement}

Versioned contracts and repository interfaces allow later support for additional corridors, several emergency vehicles, new observation sources, distributed training, or a different simulator. These are extension points rather than current capabilities. Any extension must preserve single ownership, causal event ordering, legal phase transitions, and explicit comparison manifests so increased scale does not weaken the evidence model. Multiple emergency vehicles would require a conflict-resolution policy rather than simply creating more corridor requests. A live camera source would require privacy review, calibration, clock synchronization, and monitored failure behavior. Distribution would require authenticated service boundaries and a stronger event-delivery model, while a simulator replacement would require an adapter capable of reproducing the same snapshot and applied-action semantics. Keeping these possibilities outside the current claim avoids implying unfinished functionality, yet documenting their architectural seams prevents short-term choices from closing a reasonable development path.

\subsection{Version Control}

Git maintains source code, schemas, scenarios, and report figures. Each retained experiment records its commit identifier, dependency lock, model checksum, network and demand versions, controller configuration, and metric schema. Demonstration and result releases are tagged. Generated outputs do not replace source configuration, and secrets or machine-specific paths are excluded from portable manifests. Feature work is reviewed through branches or equivalent isolated changes before it is incorporated into an evidence-producing release. Schema and contract changes require coordinated fixtures and migration notes; a commit that changes the meaning of a metric also increments its definition version. Large model, video, and run artifacts are referenced by checksums rather than committed indiscriminately. This policy links every reported result to an inspectable software state while keeping the repository manageable. It also allows an evaluator to distinguish an exact reproduction from a later rerun made after code, data, or configuration changed.

\subsection{User Interface Paradigm}

The interface uses task-oriented screens and a unidirectional command-and-view model. Operators configure runs, monitor network and mission state, inspect provenance, and review evidence without directly manipulating the simulator. Validation errors are tied to fields; colors have textual labels; units and freshness remain visible; and measured results remain distinct from schematic values. These rules implement Section \ref{sec:gui} and permit browser reconnection without transferring control ownership. Screen identifiers and navigation reflect the user's workflow from configuration through monitoring to comparison, while specialist details remain available without crowding the overview. Emergency state is displayed with route version, lifecycle state, and warnings instead of a single ambiguous priority indicator. Analytics disclose sample count, variability, failed runs, metric definition, and configuration differences near each comparison. The interface therefore supports operational awareness and academic review, but it never presents unimplemented fields, schematic charts, or stale streams as measured evidence.

\subsection{External Database and Data Storage Management}

Structured configuration, events, routes, and metrics follow Section \ref{sec:data-design}; videos, models, and logs remain checksummed external artifacts. A repository interface permits file-backed integration and later PostgreSQL use without changing domain meaning. Schema migrations, referential constraints, immutable completed-run configuration, and versioned metrics protect comparison integrity. Bounded buffering may tolerate a temporary write failure, but processing stops before evidence loss becomes unbounded. The database records identifiers and provenance for external artifacts so a run does not depend on an unexplained local filename. Observation ownership is exclusive, route versions increase monotonically, and priority events retain the route that justified them. Applied control records are authoritative for evaluation because a safety intervention may differ from the controller request. Backup and export procedures include manifests and schema versions, while credentials and licensed content remain outside research exports. These measures keep storage decisions aligned with auditability rather than treating persistence as a passive logging concern.

\subsection{Concurrency and Synchronization}

The orchestrator is the only simulation-step owner, and TraCI access is serialized through the adapter. Controllers receive immutable snapshots and return requests instead of mutating shared state. Perception uses a bounded queue, dashboard streams may skip views without delaying control, and route publication is atomic with cancellation of obsolete corridor work. Shutdown rejects new commands, terminates or detaches SUMO safely, flushes records, and assigns an explicit terminal state. Simulation time, monotonic event sequence, and wall-clock time are stored separately because inference and persistence delays must not be mistaken for traffic causality. A late observation is marked stale or omitted according to policy instead of blocking the loop indefinitely. The active route version changes together with invalidation of its old requests, preventing two paths from owning the same corridor. Database batches may improve throughput, but applied actions and mission transitions must first enter a durable bounded queue. These rules define synchronization even when all components run on one machine.

\section{Domain Model/Class Diagram}
\label{sec:domain-model}

\begin{figure}[htbp]
\centering
\begin{tikzpicture}[font=\footnotesize,node distance=8mm and 8mm,cls/.style={draw,rounded corners,fill=manzilgray,minimum width=3.15cm,minimum height=1.05cm,align=left},assoc/.style={-{Latex},thin}]
\node[cls] (scenario) {\textbf{Scenario}\\network, demand};
\node[cls,right=of scenario] (run) {\textbf{ExperimentRun}\\id, seed, status};
\node[cls,right=of run] (snap) {\textbf{TrafficSnapshot}\\source, step, validity};
\node[cls,below=of scenario] (mission) {\textbf{EmergencyMission}\\origin, destination, status};
\node[cls,below=of run] (orch) {\textbf{Orchestrator}\\step(), schedule()};
\node[cls,below=of snap] (decision) {\textbf{ControlDecision}\\requested, applied, reason};
\node[cls,below=of mission] (route) {\textbf{RouteVersion}\\edges, cost, trigger};
\node[cls,below=of orch] (ctrl) {\textbf{Controller}\\requestAction()};
\node[cls,below=of decision] (safety) {\textbf{SafetyFilter}\\validate(), transform()};
\node[cls,below=of route] (corr) {\textbf{CorridorState}\\intersection, state};
\node[cls,below=of safety] (adapter) {\textbf{SimulationAdapter}\\advance(), apply()};
\draw[assoc] (scenario)--node[midway,above=2mm,font=\scriptsize]{1:N}(run);
\draw[assoc] (run)--node[right,font=\scriptsize]{1}(orch); \draw[assoc] (orch)--(ctrl);
\draw[assoc] (run)--node[midway,above,sloped,font=\scriptsize]{0..1 active}(mission);
\draw[assoc] (mission)--node[left,font=\scriptsize]{1:N}(route); \draw[assoc] (route)--node[left,font=\scriptsize]{0:N}(corr);
\draw[assoc] (snap)--(decision); \draw[assoc] (ctrl)--(decision); \draw[assoc] (decision)--(safety); \draw[assoc] (corr)--(safety); \draw[assoc] (safety)--(adapter);
\end{tikzpicture}
\caption{Proposed domain model/class diagram}
\label{fig:class}
\end{figure}

Figure \ref{fig:class} presents the core domain objects and their relationships independently of a particular web or database framework. A Scenario can produce many ExperimentRun records, while each run has one orchestrator that schedules controllers and the sole simulation adapter. The model deliberately separates a ControlDecision's requested phase from its applied phase. A controller can request an early or otherwise invalid switch, but the SafetyFilter validates or transforms that request before the adapter can reach SUMO. TrafficSnapshot carries source, simulation step, and validity so a decision can be traced to the observation it used. A run has at most one active EmergencyMission in the core scope. That mission owns immutable RouteVersion objects, and each version can justify several CorridorState records. Connecting the corridor to the same SafetyFilter makes clear that normal and emergency actions share legal-transition enforcement. Exact method signatures may change during implementation, but the cardinalities, single-writer path, route versioning, and requested-versus-applied distinction must remain intact.

\section{Sequence Diagrams}
\label{sec:sequences}

\subsection{Normal Control Update}

\begin{figure}[htbp]
\centering
\begin{tikzpicture}[font=\footnotesize,x=2.2cm,y=.62cm]
\foreach \x/\n in {0/Orchestrator,1/Adapter,2/State,3/Controller,4/Experiment,5/Store}{\node at (\x,0) {\n};\draw[dashed] (\x,-.3)--(\x,-8.2);}
\draw[-{Latex}] (0,-1)--node[above]{advance(step)}(1,-1);
\draw[-{Latex}] (1,-2)--node[above]{raw observations}(2,-2);
\draw[-{Latex}] (2,-3)--node[above]{snapshot(t)}(0,-3);
\draw[-{Latex}] (0,-4)--node[above]{request(snapshot)}(3,-4);
\draw[-{Latex}] (3,-5)--node[above]{phase request}(0,-5);
\draw[-{Latex}] (0,-6)--node[above]{validated action}(1,-6);
\draw[-{Latex}] (1,-7)--node[above]{applied event}(4,-7);
\draw[-{Latex}] (4,-8)--node[above]{validated record}(5,-8);
\end{tikzpicture}
\caption{Sequence for a normal signal-control update}
\label{fig:seq-normal}
\end{figure}

Figure \ref{fig:seq-normal} shows the causal order for one normal-control decision. The orchestrator first asks the adapter to advance one simulation step, after which raw subscriptions are normalized into a complete snapshot. Only then does the selected controller receive state and return a phase request. The orchestrator validates that request against phase membership, elapsed minimum green, clearance, and ownership before the adapter applies it. The resulting event passes through the experiment service for schema validation and run association before the repository stores it. A controller therefore never observes a partial update, cannot actuate SUMO as a side effect, and cannot write an unvalidated evidence record. When a controller decision interval spans several simulation steps, the last applied action remains in force while state and metrics continue to be sampled at their declared intervals. The stored event includes requested and applied action, effective step, snapshot step, and intervention reason, allowing replay to distinguish a policy choice from a safety transformation.

\subsection{Ambulance Dispatch and Corridor Activation}

\begin{figure}[htbp]
\centering
\begin{tikzpicture}[font=\footnotesize,x=1.9cm,y=.60cm]
\foreach \x/\n in {0/Dashboard,1/API,2/Routing,3/Corridor,4/Adapter,5/Experiment,6/Store}{\node at (\x,0) {\n};\draw[dashed] (\x,-.3)--(\x,-9.2);}
\draw[-{Latex}] (0,-1)--node[above]{dispatch(o,d)}(1,-1);
\draw[-{Latex}] (1,-2)--node[above]{validate + compute}(2,-2);
\draw[-{Latex}] (2,-3)--node[above]{route v1}(1,-3);
\draw[-{Latex}] (2,-4)--node[above]{look-ahead requests}(3,-4);
\draw[-{Latex}] (3,-5)--node[above]{safe transition request}(4,-5);
\draw[-{Latex}] (4,-6)--node[above]{applied state}(3,-6);
\draw[-{Latex}] (3,-7)--node[above]{priority events}(5,-7);
\draw[-{Latex}] (5,-8)--node[above]{validated records}(6,-8);
\draw[-{Latex}] (1,-9)--node[above]{mission view}(0,-9);
\end{tikzpicture}
\caption{Sequence for ambulance dispatch and corridor activation}
\label{fig:seq-dispatch}
\end{figure}

Figure \ref{fig:seq-dispatch} begins when the operator submits an origin and destination through the dashboard. FastAPI validates the command and asks the routing service to resolve accessible edges and compute the first immutable route version. A successful route is returned for the mission view and supplied to the corridor coordinator, which creates bounded look-ahead requests rather than activating the whole path. Each due request passes through safe transition before the adapter applies priority. The corridor sends mission and priority events to the experiment service, which validates their identifiers and route versions before persistence. Invalid endpoints or a disconnected graph stop the sequence before an ambulance or partial corridor is activated. The dashboard receives a derived mission view after authoritative state has changed, so disconnecting the client cannot cancel or delay the mission. The sequence separates route computation, signal ownership, evidence validation, and presentation while retaining the identifiers required to trace an active priority state back to the command and route version.

\subsection{Rerouting and Recovery}

\begin{figure}[htbp]
\centering
\begin{tikzpicture}[font=\footnotesize,x=2.2cm,y=.60cm]
\foreach \x/\n in {0/State,1/Routing,2/Corridor,3/Adapter,4/Experiment,5/Store}{\node at (\x,0) {\n};\draw[dashed] (\x,-.3)--(\x,-9.0);}
\draw[-{Latex}] (0,-1)--node[above]{cost/closure change}(1,-1);
\draw[-{Latex}] (1,-2)--node[above]{candidate + threshold}(1,-2.8);
\draw[-{Latex}] (1,-3.5)--node[above]{route v2}(2,-3.5);
\draw[-{Latex}] (2,-4.5)--node[above]{cancel obsolete}(3,-4.5);
\draw[-{Latex}] (2,-5.5)--node[above]{request new look-ahead}(3,-5.5);
\draw[-{Latex}] (0,-6.5)--node[above]{passage confirmed}(2,-6.5);
\draw[-{Latex}] (2,-7.5)--node[above]{recovery action}(3,-7.5);
\draw[-{Latex}] (2,-8.3)--node[above]{revision/events}(4,-8.3);
\draw[-{Latex}] (4,-8.9)--node[above]{validated records}(5,-8.9);
\end{tikzpicture}
\caption{Sequence for route revision, obsolete-priority cancellation, and recovery}
\label{fig:seq-reroute}
\end{figure}

Figure \ref{fig:seq-reroute} shows how a closure or sustained cost change is handled after dispatch. The routing service evaluates a candidate from the ambulance's valid current position and applies the configured feasibility, benefit, persistence, and cooldown rules. If accepted, a new immutable route version is published to the corridor coordinator. Requests associated with the old remaining path are cancelled before the new look-ahead set is activated, preventing simultaneous priority for incompatible route versions. Passage confirmation still originates from traffic state and moves the relevant intersection into recovery, while route revisions, cancellations, requests, and recovery events are stored in causal order. When no feasible replacement exists, the mission becomes blocked and upcoming priority is cancelled rather than preserving a corridor that no longer reaches the destination. This interaction makes rerouting an auditable state transition, not merely a line redrawn on the dashboard. Rejected candidates also remain available for diagnosis.

\section{Policies and Tactics}

\subsection{Corridor State and Safety Policy}
\label{sec:corridor-state}

\begin{figure}[htbp]
\centering
\begin{tikzpicture}[font=\footnotesize,node distance=14mm and 12mm,st/.style={draw,rounded corners,minimum width=2.45cm,minimum height=.85cm,align=center,fill=manzilgray},arr/.style={-{Latex},thick}]
\node[st] (idle) {IDLE};
\node[st,right=of idle] (req) {REQUESTED};
\node[st,right=of req,fill=manzilamber!25] (trans) {SAFE TRANSITION};
\node[st,right=of trans,fill=manzilgreen!20] (active) {ACTIVE PRIORITY};
\node[st,below=of active] (passed) {PASSAGE};
\node[st,left=of passed] (recover) {RECOVERY};
\node[st,left=of recover,fill=manzilred!12] (cancel) {CANCEL/FAILURE};
\draw[arr] (idle)--node[above,font=\scriptsize]{valid}(req);
\draw[arr] (req)--node[above,font=\scriptsize]{due}(trans);
\draw[arr] (trans)--node[above,font=\scriptsize]{cleared}(active);
\draw[arr] (active)--node[right,font=\scriptsize]{passes}(passed);
\draw[arr] (passed)--(recover);\draw[arr] (recover.south) -- ++(0,-6mm) -| node[pos=.45,below,font=\scriptsize]{normal boundary} (idle.south);
\draw[arr] (req) to[bend right=42] (idle);
\draw[arr] (trans.south) -- node[right,font=\scriptsize]{cancel/fail} (cancel.north);
\draw[arr] (active.west) -| (cancel.east);
\draw[arr] (cancel)--node[below,font=\scriptsize]{safe fallback}(recover);
\end{tikzpicture}
\caption{Intersection-level moving-corridor state transitions}
\label{fig:corridor-state}
\end{figure}

Figure \ref{fig:corridor-state} defines the lifecycle owned by each affected intersection. IDLE accepts a valid route-versioned request, but REQUESTED does not imply signal ownership. A request that expires or becomes invalid at this stage returns directly to IDLE because no signal state has changed and recovery would be misleading. When activation becomes due, SAFE TRANSITION waits for the current phase's minimum service and completes the configured yellow and all-red clearance. ACTIVE PRIORITY serves only the ambulance movement and is bounded by a maximum duration. Confirmed clearance creates a PASSAGE event, after which RECOVERY restores a valid normal boundary before ordinary control resumes. Cancellation, timeout, or route revision during SAFE TRANSITION or ACTIVE PRIORITY enters the guarded failure and recovery path because signal ownership may already have changed. The distinction prevents an unowned queued request from triggering unnecessary compensation while ensuring that every applied intervention returns safely to IDLE. Each transition records its trigger, route version, and effective step.

\subsection{Interface and Event Contracts}

Table \ref{tab:contracts} lists the minimum fields carried across important module boundaries. Version and time fields prevent a consumer from accepting an unknown schema or treating an old snapshot as current. Source mode, units, confidence, and validity distinguish simulator values from calibrated or proxy video observations. SignalState identifies both the legal successors and current owner, while ControlRequest has no direct side effect and can be transformed into an authoritative AppliedControl record by the safety path. RoutePlan and PriorityRequest carry the mission and route version needed to reject obsolete corridor work. MetricSample retains its definition version and interval so summaries cannot combine measurements whose meaning changed during development. Field names are conceptual and may be mapped to Python or JSON naming conventions, but their semantics remain stable. This contract-first approach also makes integration tests possible before every research component is complete, because a stub can produce valid messages without bypassing ownership rules.

\begin{longtable}{p{0.20\textwidth}p{0.21\textwidth}p{0.46\textwidth}}
\caption{Core interface and data-contract fields}\label{tab:contracts}\\
\toprule
\textbf{Contract} & \textbf{Required fields} & \textbf{Rule}\\
\midrule\endfirsthead
\toprule
\textbf{Contract} & \textbf{Required fields} & \textbf{Rule}\\
\midrule\endhead
TrafficSnapshot & schemaVersion, sourceMode, producedAt, effectiveStep, validUntil, networkId & Consumers reject unknown schema or expired validity; wall time and simulation time are distinct.\\
ApproachState & intersectionId, approachId, count, queue, wait, speed, occupancy, units, confidence & Missing is not zero; uncalibrated physical quantities are omitted.\\
SignalState & intersectionId, programId, phaseId, phaseElapsed, legalNext, owner & Owner is NORMAL or CORRIDOR; legal next phases come from validated configuration.\\
ControlRequest & controllerId, intersectionId, requestedPhase, duration/extension, snapshotStep, reason & Request has no side effect and may be rejected or transformed.\\
AppliedControl & requestId, appliedPhase, effectiveStep, safetyIntervention, reason & This record is authoritative for analysis.\\
RoutePlan & missionId, routeVersion, edgeIds, costSeconds, snapshotStep, trigger & Edge order must be continuous from valid current position.\\
PriorityRequest & missionId, routeVersion, intersectionId, movementId, ETA, expiresAt & Obsolete route version or expired request is rejected.\\
MetricSample & runId, metricName, metricVersion, interval, value, unit & Aggregation across mismatched versions is prohibited.\\
\bottomrule
\end{longtable}

Event ordering is based on simulation step within a run and a monotonic sequence number assigned by the orchestrator. Wall-clock timestamps support diagnostics but do not determine traffic causality because inference, rendering, and storage latency can vary independently of simulated time. A snapshot identifies the step at which its values became effective, and an applied signal event records the step at which the adapter enacted it. Route revisions and priority events are enqueued in an order that prevents an old request from being interpreted as current. Missing values remain missing rather than being converted to zero, and unknown schema versions are rejected rather than partially decoded. These rules give the fields in Table \ref{tab:contracts} operational meaning and support replay, stale-data handling, and comparison across runs without assuming perfectly synchronized wall clocks. They also allow a delayed write to be distinguished from a delayed decision during failure analysis and experimental replay.

\subsection{Resource Ownership Tactics}

The orchestrator is the only writer to simulation time, and TraCI access is serialized behind the adapter. Controller computation receives an immutable snapshot and returns a request rather than holding a reference to mutable global state. Perception uses a bounded work queue; if inference falls behind, the worker marks observations late or drops work according to a recorded policy instead of blocking control indefinitely. Dashboard streams are lossy views, so a slow client may skip intermediate snapshots but cannot delay decisions. Run configuration becomes immutable after start. The current route version is updated atomically with cancellation of obsolete corridor requests, preventing new and old paths from owning priority concurrently. Database writes may be batched where safe, but an applied signal event or mission transition is not acknowledged until it has entered a durable bounded queue. During shutdown, the API stops accepting new commands, the orchestrator ends or detaches SUMO safely, queued records are flushed, and the run receives an explicit completed, stopped, failed, or incomplete terminal state.

\subsection{Error Handling and Recovery}

Table \ref{tab:errors} defines immediate containment and retained evidence for failures at each major boundary. A SUMO or TraCI failure stops decision-making because continuing to show live state would be misleading, whereas a dashboard disconnect leaves the authoritative experiment unaffected. Stale observations and invalid controller requests are rejected or handled through a declared simulator or fixed-control fallback, and every interval using that fallback is marked so it cannot be included silently in a clean learned-controller aggregate. A missing route prevents partial corridor activation, while a lost ambulance or timeout cancels future priority and moves active intersections through recovery. Persistence uses a bounded buffer only long enough to tolerate a transient issue; the run stops before evidence loss grows without limit and is marked incomplete if durability cannot be restored. These policies favor an explicit failed run over an apparently successful record whose causal events are missing. Fault-injection and small-network integration tests will verify both the immediate behavior and the final recorded state.

\begin{table}[htbp]
\centering
\caption{Failure handling strategy}
\label{tab:errors}
\begin{tabularx}{\textwidth}{p{0.27\textwidth}YY}
\toprule
\textbf{Failure} & \textbf{Immediate behavior} & \textbf{Recorded/recovery behavior}\\
\midrule
SUMO launch/TraCI disconnect & Stop decisions and do not present live state & Mark run failed; retain command, stderr, step, and manifest\\
Stale observation & Exclude or use declared simulator fallback & Record staleness and selected fallback source\\
Invalid controller action & Reject through safety filter & Apply fixed/safe continuation and record intervention\\
RL service unavailable & Switch at valid boundary to configured baseline & Record fallback interval; exclude from clean RL aggregate or report separately\\
No route & Do not spawn/activate partial corridor & Mark mission blocked; show cause and graph snapshot ID\\
Lost ambulance/timeout & Cancel future priority and recover active intersections & Mark mission failed/cancelled and preserve last position\\
Persistence failure & Buffer within a bounded queue; stop before unbounded evidence loss & Mark run incomplete if durability cannot be restored\\
Dashboard disconnect & Continue experiment unaffected & Reconnect from current snapshot and event cursor\\
\bottomrule
\end{tabularx}
\end{table}

\subsection{Implementation Policies}

\subsubsection{Controller Policy}

All controllers implement a common method that accepts an immutable traffic snapshot, intersection configuration, and decision context, then returns a request with explanation fields. They cannot write global state, call TraCI, or suppress the record of a rejected action. The fixed controller, pressure-inspired controller, and learned policy therefore differ in decision logic while sharing phase identifiers, decision intervals, safety constraints, and applied-event recording. Deterministic evaluation disables exploration and seeds every permitted stochastic operation. The safety filter verifies phase membership, minimum green, required clearance, maximum priority duration, and current ownership before an action reaches the adapter. If a policy is unavailable or repeatedly invalid, control transfers at a legal boundary to the configured baseline and the affected interval is marked. This common interface is central to a fair comparison because it prevents one controller from receiving more permissive movements or direct simulator privileges. Explanation fields are retained for later error analysis.

\subsubsection{Routing Policy}

The graph builder derives nodes and edges from validated SUMO vehicle movements, including allowed classes and turn continuity. Cost providers are composable but must state units, observation time, fallback behavior, and whether link travel time already includes queue delay. This metadata prevents the routing service from adding overlapping estimates or comparing seconds with an unnormalized congestion score. A published route plan is immutable, and the mission points to exactly one current version. Revisions are accepted only through Algorithm \ref{alg:routing}, which compares paths under one cost snapshot and applies feasibility, hysteresis, persistence, and cooldown rules. A hard closure overrides ordinary cost and invalidates an affected remaining path immediately. Before publication, the route is checked from the ambulance's current edge to the destination and given a revision cause. The routing service never actuates signals; it supplies versioned plans to the corridor coordinator, which can then cancel or create priority requests consistently.

\subsubsection{Persistence Policy}

The store uses schema migrations and referential constraints for structured configuration, events, missions, routes, and metrics, while large videos, model files, and logs may remain checksummed external artifacts with provenance. Raw high-frequency state can be downsampled for persistence only when the sampling interval and aggregation rule are recorded in the metric definition. Completed run configuration is immutable, and derived summaries retain their source run identifiers and analysis version. Secrets, machine-specific credentials, and unportable local paths are excluded from research exports; licensed artifacts are referenced rather than redistributed when necessary. PostgreSQL is proposed for structured records, but implementation may begin with a file-backed engine behind the same repository interface. The domain schema and integrity rules remain independent of that choice so a storage migration cannot change the meaning of an experiment. A persistence failure follows the bounded behavior in Table \ref{tab:errors} rather than silently discarding events. Recovery records state whether buffered evidence was restored.

\subsubsection{User-Interface Policy}

The user interface is a view and command boundary, not a second control plane. Colors are supplemented with text and symbols so congestion, corridor state, or failure is not conveyed by color alone. Values display their units, source mode, freshness, and invalid state where applicable. Schematic predictions remain visually distinct from measured results, and an uncalibrated video proxy is not labeled as a physical queue or speed. The interface does not display ``time saved'' until a valid reference run exists under the same comparison manifest, and it exposes configuration differences when direct comparison is unsafe. Commands receive validation errors tied to fields, while connection loss leaves the experiment running and allows a client to reconnect from the current event cursor. These policies implement the wireframe intentions in Section \ref{sec:gui} and ensure that a polished display cannot conceal stale data, fallback intervals, missing runs, or incomplete evidence.

\subsection{Testing Strategy}

Unit tests cover cost functions, route continuity, hysteresis, metric calculations, schema validation, freshness rules, and corridor transition guards. Property or invariant tests generate action sequences to ensure incompatible greens are never requested for application, one component owns each intersection, obsolete route versions cannot activate priority, and cancellation eventually reaches IDLE through RECOVERY. Integration tests use small SUMO networks to validate single-step ownership, detector-to-state messages, route application, fallback, and persistence ordering. Scenario tests replay fixed seeds for normal control, closure rerouting, no-route, invalid actions, stale data, controller loss, ambulance timeout, persistence failure, and dashboard disconnect. Contract fixtures allow components to be replaced by stubs without weakening field validation. Research verification remains separate from software correctness: passing these tests shows that the implementation follows its specification, but it does not establish that a policy improves traffic or emergency travel. The matched experiments in Chapter \ref{ch:method} provide that evidence, while test results explain whether a failed experiment reflects the policy or a software defect.

\subsection{Coding, Build, and Maintainability Policies}

Python modules will use static type hints, consistent formatting and linting, bounded dependency versions, structured logging, and documentation for public interfaces. React components separate data acquisition from presentation so a visual revision cannot change experiment ownership. Configuration is stored in validated versioned files rather than scattered constants, and secrets remain outside exported manifests. Continuous checks should run unit tests, schema validation, and a minimal headless SUMO scenario before a revision is used for evidence. Experiment releases are tagged and their commit identifier is written into each run manifest. The intended repository structure separates application services, controllers, routing, perception, simulation, dashboard, contracts, scenarios, and experiment artifacts. This is a proposed organization, not evidence that all directories or services already exist. Small interfaces allow a tracker, learned policy, database engine, or dashboard transport to change without rewriting the central loop, while versioned contracts force such changes to declare any effect on recorded evidence.

\subsection{Continuous Integration and Deployment}

The proposed continuous-integration pipeline runs formatting, linting, type checks, unit tests, contract fixtures, schema validation, and a minimal headless SUMO scenario. A build that changes an interface or metric definition must update its version and associated fixtures. Deployment for the Deliverable 2 prototype remains a documented local environment rather than a claimed production service. A tagged demonstration build includes dependency locks, validated scenarios, required model checksums, and migration state so another team member can reproduce the same prepared workflow. The headless smoke scenario verifies process launch, one controlled step sequence, valid signal application, event persistence, and a clean terminal state without requiring the dashboard. Research experiments are not automatically promoted merely because software tests pass; the matched evaluation protocol remains a separate activity. Failed checks block an evidence release, and bypasses must be recorded rather than hidden. If an external model or licensed video cannot be packaged, the build manifest records its checksum and acquisition instructions while keeping the protected artifact outside the archive.

\subsection{Source Code Organization}

The intended repository separates API and application services from controllers, routing, perception, simulation adapters, interface contracts, dashboard code, scenarios, and experiment artifacts. Tests mirror those boundaries, while generated outputs and large external artifacts remain outside source modules. Configuration is versioned and validated instead of embedded in code. This arrangement keeps research implementations replaceable, makes ownership visible during review, and avoids circular dependencies between interface, simulation, and persistence concerns. Domain objects and contracts occupy a low-dependency package that does not import the dashboard or infrastructure adapters. Controller implementations share an interface but do not import each other, allowing fixed, pressure-inspired, and learned policies to be selected by configuration. Scenario manifests reference networks and demand files through portable paths and checksums. Analysis scripts read retained records rather than private controller state. Documentation accompanies public interfaces and configuration schemas. This organization remains proposed until source implementation confirms it, but it defines a reviewable target that aligns folders, dependency direction, component ownership, and the architectural diagrams in this chapter.

\section{Conclusion}

The design assigns one clear owner to every mutable resource and makes observation source, simulation time, units, route version, requested action, and applied action explicit. The high-level architecture keeps presentation and perception outside the single-writer simulation path, while component responsibilities and data contracts make integration boundaries testable. The domain model and sequences preserve causal links from scenario and snapshot through controller request, safety decision, route revision, corridor transition, and durable event. Guarded recovery and bounded fallbacks ensure that a failure does not leave an unexplained signal or apparently live dashboard. Common controller interfaces support the intended progression from fixed and pressure-inspired baselines to a learned policy, and versioned routing enables a moving corridor without granting the pathfinder direct signal authority. The architecture is modular enough for team implementation on one machine and leaves distribution as a profiling decision, not an assumption. It also preserves failed runs and negative research outcomes as evidence instead of hiding them behind interface behavior.
